% Generated draft: requires mathematical and source review.
\documentclass[11pt]{article}
\usepackage[margin=25mm]{geometry}
\usepackage{fontspec}
\setmainfont{Latin Modern Roman}
\usepackage{amsmath,amssymb,mathtools,unicode-math}
\setmathfont{Latin Modern Math}
\usepackage{enumitem,needspace,xurl,hyperref}
\setlength{\parindent}{0pt}
\setlength{\parskip}{5pt}
\setlength{\emergencystretch}{3em}
\newcommand{\N}{\mathbb{N}}
\newcommand{\Z}{\mathbb{Z}}
\newcommand{\Q}{\mathbb{Q}}
\newcommand{\R}{\mathbb{R}}
\newcommand{\C}{\mathbb{C}}
\newcommand{\F}{\mathbb{F}}
\newcommand{\A}{\mathbb{A}}
\newcommand{\PP}{\mathbb{P}}
\newcommand{\E}{\mathbb{E}}
\newcommand{\eps}{\varepsilon}
\newcommand{\dd}{\,\mathrm{d}}
\begin{document}
\section*{UnsolvedMath ID: 701081}
\textbf{Source problem code:} catalogue entry
\subsection{Definitions and mathematical statement}
For a group \(G\), write \(|g|\) for the order of an element \(g\in G\). The group is periodic if every element has finite order. For distinct primes \(p,q\), a periodic \(\{p,q\}\)-group is a group in which every element has order \(p^a q^b\) for some integers \(a,b\geq 0\), depending on the element. A periodic group is primary if there is a single prime \(r\) such that every element has order a power of \(r\). Thus, in a non-primary periodic \(\{p,q\}\)-group, elements of order \(p\) and elements of order \(q\) both occur.

A group \(G\) is locally soluble if every subgroup generated by a finite subset of \(G\) is soluble. Here a group \(H\) is soluble if its derived series, defined by \(H^{(0)}=H\) and \(H^{(n+1)}=[H^{(n)},H^{(n)}]\), reaches the trivial subgroup after finitely many steps. The required number of steps may depend on the finitely generated subgroup. A \(p\)-subgroup means a subgroup all of whose elements have finite order a power of \(p\); it need not be finite.

The question is whether, for every pair of distinct primes \(p,q\), there exist a non-primary periodic locally soluble \(\{p,q\}\)-group \(G\) and subgroups \(A,B\leq G\) such that both \(A\) and \(B\) are \(p\)-subgroups and
\[
G=AB=\{ab:a\in A,\ b\in B\}.
\]
The equality is an equality of sets: every element of \(G\) must be a product of one element of \(A\) followed by one element of \(B\). Neither normality nor commutativity of the factors is assumed. No finiteness assumption is imposed on \(G\), \(A\), or \(B\).
\subsection{Short English statement}
For every two distinct primes \(p\) and \(q\), does there exist a periodic locally soluble group whose element orders involve only these primes, with both primes occurring, that is the setwise product of two \(p\)-subgroups?
\subsection{Sources}
\begin{itemize}
\item Local UnsolvedMath registry, permanent ID \(701081\), source problem code \texttt{catalogue entry}, title \emph{Kourovka Notebook Problem 14.97}. The supplied record labels the problem open; that label is supplied curation, not an independently verified mathematical assessment.
\item N. S. Chernikov, Problem 14.97, \emph{The Kourovka Notebook (Unsolved Problems in Group Theory)}, 14th issue (1999), as attributed in the supplied record.
\item The supplied background also identifies the 21st issue edition (2026), \texttt{arXiv:1401.0300v48}. The collection description names E. I. Khukhro and V. D. Mazurov as editors. These bibliographic details are reproduced from the supplied record without independent source verification.
\end{itemize}
\end{document}
